\documentclass[notitlepage,onecolumn,showpacs,11pt]{revtex4-1}
\usepackage{amssymb}
\usepackage{graphicx}
\usepackage{amsmath}
\usepackage{float}
\usepackage{epstopdf}
\usepackage{times}
\usepackage{color}
\usepackage{subfigure}
\usepackage{setspace}
\usepackage{bm}% bold math

\begin{document}

\title{Disordered $J-J'$ Model in a Triangular Lattice; A Mean-field Approximation}

\maketitle

The disordered next-nearest neighbour Heisenberg model is described by the 
Hamiltonian, 
\begin{equation}
	\begin{aligned}
		H=\sum_{\langle i,j\rangle}J_{ij}\boldsymbol{S}_{i}.\boldsymbol{S}_{j}+
		\sum_{\langle\langle i,j\rangle\rangle}J'_{ij}\boldsymbol{S}_{i}.
		\boldsymbol{S}_{j}
	\end{aligned}
\end{equation}
, where both the nearest and the next nearest coupling, $ J $ and $ J'$ are 
drawn from random positive definite distributions. The distribution is taken 
to be positive definite so that the system remains anti-ferromagnetic.
To find a mean-field description of this model it is useful to trade 
the spin operators for their slave fermion representation which makes 
way for mean-field decoupling in various channels of excitation. The 
generalised SU(N) model can be exactly solved in this paradigm. However, 
such a scheme fails to capture any emergent magnetic order of the system 
and therefore remains ill-suited to describe the magnetic phases of the 
disordered system. An Ising type mean-field description which approximates 
the model in terms of local moments and local fields remain an avenue to 
describe the magnetic phases. In the following we concentrate on a hybrid 
scheme which involves both these types of mean-field decoupling by introducing 
a cosmetic parameter.

\section{The Mean-field decoupling}

The description of the mean-field decoupling are the same for both the nearest
and the next-nearest neighbour part of the hamiltonian. In the slave fermion
representation for the $ SU(2) $ model, the spin operators can be expressed as
$\bm{S}_i=f_{i\alpha}^\dagger\bm{\sigma}_{\alpha\beta}f_{i\beta}$ with the 
constraint $f^\dagger_{i\alpha}f_{i\alpha}=1$. In the notation that follows, 
only the fermion flavours (greek letters) are summed over if repeated and
the sum over the lattice co-ordinates are explicitly mentioned. An additional
parameter $x$ allows us to include both the decouping in the fermion-fermion
channel and also on the basis of the local moments. The mean-field hamiltonian 
following this scheme is obtained to be,
\begin{equation}
	\begin{aligned}
		H_{\text{mf}}&=-\sum_{\langle i,j\rangle}
		\left(\chi^{N*}_{i,j}
			f^\dagger_{i\alpha}f_{j\alpha}
		+\text{h.c.}\right)-\sum_{\langle\langle i,j\rangle\rangle}
		\left(\chi^{NN*}_{ij}
			f^\dagger_{i\alpha}f_{j\alpha}
		+\text{h.c.}\right)\\ 
		&+\frac{1}{2}\sum_{i}\boldsymbol{B}_i.f_{i\alpha}^\dagger\bm{\sigma}_{\alpha\beta}
		f_{i\beta}
		+\sum_i\lambda_i\left(f^\dagger_{i\alpha}f_{i\alpha}-1\right)
	\end{aligned}
\end{equation}
with the following set of equation determining the mean-field parameters.
\begin{equation}
	\begin{aligned}
		\langle f^\dagger_{i\alpha}f_{i\alpha}\rangle &= 1 \\
		\chi_{ij}^{N}&=(1-x)J_{ij}\langle f^\dagger_{i\alpha}f_{j\alpha}\rangle
		\text{,\ \ i,j \ }\epsilon\text{ \ N}   \\
		\chi_{ij}^{N}&=(1-y)J_{ij}\langle f^\dagger_{i\alpha}f_{j\alpha}\rangle 
		\text{,\ \ i,j \ }\epsilon\text{ \ NN}   \\
		\boldsymbol{B}_i&=\frac{1}{2}
		\left[ x\sum_{j\epsilon N}
			J_{ij}\langle f_{j\alpha}^\dagger\bm{\sigma}_{\alpha\beta}f_{j\beta}\rangle
			+ y\sum_{j \epsilon NN}
		J_{ij}\langle f_{j\alpha}^\dagger\bm{\sigma}_{\alpha\beta}f_{j\beta}\rangle\right]
	\end{aligned}
\end{equation}
It is possible to solve this system of equation in real space using exact diagonalisation 
and iterative convergence procedure.

\section{Hubbard-Stratanovich Transformation: Formal Bosonisation}

Like the mean-field decoupling expressed dexcribe it is possible to trade an interacting 
Fermionic problem for a problem with additional bosons coupled with the fermions. In its 
formulation, this constitutes an exact description of the original problem and leads way 
for further exploration using different applicable techniques. One of the formal way to 
achieve this is to perform the Hubbard-Stratanovich transformation, which, in its content 
is very much an integral identity for gaussian form. 

To illustrate the method we consider a generic interacting problem involving fermions 
$\psi_i$ in a lattice. Let the Hamiltonian be given by,
\begin{equation}\label{eq:hubbard-hamiltonian}
	\begin{aligned}
		H&=\sum_{ij}t_{ij}\psi^\dagger_i\psi_j+\sum_{ijkl}U_{ijkl}\psi^\dagger_{i}
		\psi^\dagger_{j}\psi_k\psi_l
	\end{aligned}
\end{equation}
. We are interested in the partition function of the system given by, 
$Z=\text{Tr e}^{-\beta \left(H -\mu N\right)}$ with the number operator, 
$ N=\sum_{i}\psi^\dagger_i\psi_i$. Now, to get to the functional integral formalism we 
introduce the time-slicing scheme. In the stastical partition function there is no time 
ordered integral, however, the exponent can still be written down as sum of infinite number 
of exponents, 
$\text{Tr e}^{-\beta \left(H -\mu N\right)}
=\text{Tr e}^{-\sum^{T}_{t=1} \epsilon\left(H -\mu N\right)}$ such that, $ T\epsilon=\beta $. 
This construction allows us to express the exponential in terms of multiplication of multiple
time-ordered exponentials upto some trotter error and introduce resolutions of identities 
in terms of the time-stamped Grassmann coherent states.
\subsection{Grassmann Coherent States}

Grassmann numbers anti-commute, such that for two Grassmann numbers, $\alpha\beta+\beta\alpha=0$. 
As a consequence of this anti-commutation they are also associated with the following properties,
\begin{equation}
	\begin{aligned}
		\int d\alpha 1&=0 \\
		\int d\alpha \alpha&=1 \\
		\frac{\partial}{\partial\alpha} \alpha&=0 \\
		\frac{\partial}{\partial\alpha} \alpha&=1
	\end{aligned}
\end{equation}
It is noteworthy that the integrals are defined the same way as the differentials. Unlike 
real or complex numbers these can not readily be interpreted in terms of surface areas or
infitesimal increments. However these relations are self-consistent and they unquely 
define the Grassmann calculus. At any stages of a calculation involving physical fermionic 
observables these relations can act as mnemonics and simplify the calculational approach.

Fermionic operators anti-commute so Grassmann numbers become ubiquitous for the expression 
of Grassman coherent states. The definition of the fermionic coherent states is the
following,
\begin{equation}
	\begin{aligned}
		|\psi\rangle=\text{e}^{-\sum_{i}\Psi_i\psi^\dagger_i}|0\rangle
		=\prod_i\left(1-\Psi_i\psi^\dagger_i\right)|0\rangle
	\end{aligned}
\end{equation}
, where $\Psi_i$ are Grassmann variables and $|0\rangle$ is the multi-particle vaccuam. 
The expression of the exponential in terms of multiplication here is made possible by 
the fact that the objects in the sum of the exponent commute between themselves and the 
Taylor expansion of Grassmann numbers have no non-linear non-zero terms.

It is straightforward to verify that, $\psi_i|\psi\rangle=\Psi_i|\psi\rangle $. 
Similarly for the adjoint of the coherent state we need a different flavour of Grassmann 
numbers, $\Psi^*_i$  and the expression, 
$ \langle \psi|=\langle 0|\text{e}^{-\sum_i\phi_i\Psi^*_i}$ fulfills the
second criteria as a left coherent state, $\langle\psi|\psi^\dagger=\langle\psi|\Psi^*_i$.
It can also be shown that, $\psi^\dagger_i|\psi\rangle=-\frac{\partial}{\partial\Psi_i}|\psi\rangle$
and $\langle\psi|\psi_i=+\frac{\partial}{\partial\Psi^*_i}\langle\psi|$. It is worth noting that
the different flavours of Grassmann numbers $\Psi_i$ and $\Psi^*_i$ are not related in the
usual sense of complex conjugation. The Grassmann coherent states are over-complete, 
$ \langle\psi|\psi\rangle=\text{e}^{\sum_{i}\Psi^*_i\Psi_i} $. This leads to the closure relation,
\begin{equation}
	\begin{aligned}
		\int \prod_id\Psi^*_id\Psi_i\text{e}^{-\sum_i\Psi^*_i\Psi_i}|\psi\rangle\langle\psi|&=1
	\end{aligned}
\end{equation}
Although the prefactor is apprant from the expression of over-completeness of the coherent states
it is useful to see how the above expression acts as an identity operator in the Fock space of the
many body system. We take the operator, $ A=\int \prod_id\Psi^*_id\Psi_i
\text{e}^{-\sum_i\Psi^*_i\Psi_i}|\psi\rangle\langle\psi| $ 
and find that the matrix element of this operator between two arbitrary vectors in the 
Fock space is given by,
\begin{equation}
	\begin{aligned}
		&\langle \beta_1,\beta_2,..\beta_N|A|\alpha_1,\alpha_2,..\alpha_N\rangle \\
		&=\int \prod_id\Psi^*_id\Psi_i\text{e}^{-\sum_i\Psi^*_i\Psi_i}
		\langle \beta_1,\beta_2,..\beta_N|\psi\rangle
		\langle\psi|\alpha_1,\alpha_2,..\alpha_N\rangle \\
		&=\int \prod_id\Psi^*_id\Psi_i\prod_i\left(1-\Psi^*_i\Psi_i\right)
		\langle 0|\psi_{\beta_1},\psi_{\beta_2}..\psi_{\beta_N}|\psi\rangle
		\langle\psi|\psi^\dagger_{\alpha_1},\psi^\dagger_{\alpha_2},..\psi^\dagger_{\alpha_N}|0\rangle \\
		&=\int \prod_id\Psi^*_id\Psi_i\left(1-\Psi^*_i\Psi_i\right)
		\Psi_{\beta_1}\Psi_{\beta_2}..\Psi_{\beta_N}
		\Psi^*_{\alpha_1}\Psi^*_{\alpha_2}..\Psi^*_{\alpha_N} \\ 
		&=\int \prod_id\Psi^*_id\Psi_i\left(1-\Psi^*_i\Psi_i\right)(-1)^P\delta_{\bm{\beta},\bm{\alpha}}
		\Psi_{\beta_1}\Psi_{\beta_2}..\Psi_{\beta_N}
		\Psi^*_{\beta_1}\Psi^*_{\beta_2}..\Psi^*_{\beta_N} \\
		& \ \text{for} \ P \left\{\alpha_1,\alpha_2,..\alpha_N\right\}= \left\{\beta_1,\beta_2,..\beta_N\right\}\\
		&=\int \prod_id\Psi^*_id\Psi_i\left(1-\Psi^*_i\Psi_i\right)(-1)^P\delta_{\bm{\beta},\bm{\alpha}}
		\Psi_{\beta_1}\Psi^*_{\beta_1}
		\Psi_{\beta_2}\Psi^*_{\beta_2}....\Psi_{\beta_N}\Psi^*_{\beta_N} \\
		&=(-1)^P\delta_{\bm{\beta},\bm{\alpha}}=\langle \beta_1,\beta_2,..\beta_N|\alpha_1,\alpha_2,..\alpha_N\rangle 
	\end{aligned}
\end{equation}
where the $\delta$ function determines whether the Fock index vector $ \bm{\beta}$ is same as $ \bm{\alpha}$ upto a
permutation $P$. The exponential factori $ \text{e}^{-\sum_i\Psi^*_i\Psi_i}$ is crucial for the formalism because without 
it the Grassmann integral, $ \int d\Psi^*_id\Psi_i\left(1-\Psi^*_i\Psi_i\right)\Psi_{\beta_1}\Psi^*_{\beta_1}$ would 
yield to $0$ if $ i\neq \beta_1$. So, the formalism for Grassmann integral and the Fock-space identity, 
$\int \prod_id\Psi^*_id\Psi_i\text{e}^{-\sum_i\Psi^*_i\Psi_i}|\psi\rangle\langle\psi|=1$, remains internally consistent.

\subsection{The Functional Integral Formalism and The Hubbard-Stratanovich Transformation}

With the algebra for Fermionic-coherent state set up we can return to our original problem, i.e. to describe the 
partition function of the interacting Fermion model in a lattice in terms of functional integrals. The expression for
the partition function is, $ Z=\sum_n\langle \text{n}|\text{e}^{-\beta\left(H-\mu N\right)}|\text{n}\rangle $, 
where $ |\text{n}\rangle $ are the Fock-space eigenstates of the Many-body hamiltonian. Since the overlap between 
the Fock-space eigenstates and the coherent states, $\langle \text{n}|\psi\rangle $ contain Grassmann numbers, we have the
following identity, $ \langle \text{n} |\psi\rangle\langle\psi|\text{n}\rangle=
\langle -\psi|\text{n}\rangle\langle\text{n}|\psi\rangle $. This relation is usefull to exchange the trace over
the Fock-space states with traces over Ferrmionic coherent states.
\begin{equation}
	\begin{aligned}
		Z&=\sum_n\langle\text{n}|\text{e}^{-\beta\left(H-\mu N\right)}|\text{n}\rangle \\
		 &=\sum_n\left(\int\prod_id\Psi^*_id\Psi_i\text{e}^{-\sum_i\Psi^*_i\Psi_i}\right)^2
		\langle\text{n}|-\psi\rangle\langle-\psi|\text{e}^{-\beta\left(H-\mu N\right)}|
		\psi\rangle\langle\psi|\text{n}\rangle \\
		&=\int\prod_id\Psi^*_id\Psi_i\text{e}^{-\sum_i\Psi^*_i\Psi_i}
		\langle-\psi|\text{e}^{-\beta\left(H-\mu N\right)}|\psi\rangle
	\end{aligned}
\end{equation}
Since our hamiltonian is expressed in terms
of the creation and annihilation operators, the 'time' slicing and subsequent insertion of the Fermionic coherent
states re-expresses the partion function in terms of Functional integrals over the field of time stamped Grassmann
numbers. Now, with the absence of any time integral in the exponent the formal way to introduce time slicing, 
$ N\epsilon = \beta $ would require a symbolic label for the Fermionic coherent states and their respective Grassmann fields.
\begin{equation}
	\begin{aligned}
		Z =&\int_{\psi(N)=-\psi(1)}\prod^{t=N-1 }_{i,t=1}d\Psi^*_i(t)d\Psi_i(t)
		\text{e}^{-\sum_{i,t}\Psi^*_i(t)\Psi_i(t)} \\
		&\langle\psi(N)|:\text{e}^{-\epsilon\left(H-\mu N\right)}:|\psi(N-1)\rangle..
		\langle\psi(2)|:\text{e}^{-\epsilon\left(H-\mu N\right)}:|\psi(1)\rangle
	\end{aligned}
\end{equation}
. Here, :: deontes the normal ordering of the operators (i.e. the creation operators are all to the left of the
annihilation operators) and the Trotter error arising from trancating the Baker-Campbell-Hausdroff formula 
upto the linear term (which is $\mathcal{O}(\epsilon^2)$) is ignored. The matrix-element of the exponentiated 
Hamiltonian (\ref{eq:hubbard-hamiltonian}) is given by,
\begin{equation}
	\begin{aligned}
		\langle \psi(t) |
		\text{e}^{-\epsilon\left(H-\mu N\right)}
		\psi(t-1) \rangle
		=&\langle\psi(t)|1-\epsilon
		\left(\sum_{ij}\left(t_{ij}-\mu\delta_{ij}\right)
			\psi^\dagger_i\psi_j+\sum_{ijkl}U_{ijkl}\psi^\dagger_{i}
		\psi^\dagger_{j}\psi_k\psi_l\right)|\psi(t-1)\rangle \\
		=&1+\sum_i\Psi^*_i(t)\Psi_i(t-1) \\
		 &-\epsilon
		\left(\sum_{ij}
			\left(t_{ij}-\mu\delta_{ij}\right)\Psi^*_i(t)\Psi_j(t-1)
		+\sum_{ijkl}U_{ijkl}\Psi^*_{i}(t)\Psi^*_{j}(t)\Psi_k(t-1)\Psi_l(t-1)\right) \\
		=&\text{e}^{\sum_i\Psi^*_i(t)\Psi_i(t-1)-\epsilon H(\Psi^*(t),\Psi(t-1))}
	\end{aligned}
\end{equation}
, where $ H(\Psi^*(t),\Psi(t-1)) $ is a shortend notation for the hamiltonian to remind that the two different flavours 
of Grassmann field enters in the expression of the Hamiltonian with different time-index. However, in the time-continuum 
limit when we would take,
$N\rightarrow\infty \ | \ \epsilon\rightarrow0$, we shall lose this distinction for the Hamiltonian. We substitute the 
expression for the matrix element in the original expression for the partition function and find that in the continuum
limit,
\begin{equation}
	\begin{aligned}
		Z=&\lim_{\epsilon\rightarrow 0}\int_{\psi(N)=-\psi(1)}\prod^{t=N-1 }_{t=1}
		\left(\prod_id\Psi^*_i(t)d\Psi_i(t)\right)
		\text{e}^{-\sum^{N}_{t=2}\epsilon\left[\sum_i\left(\frac{\Psi^*_i(t)\Psi_i(t)
		-\Psi^*_i(t)\Psi_i(t-1)}{\epsilon}\right)+H(\Psi^*(t),\Psi(t-1))\right]} \\
		=&\int_{\psi(\beta)=-\psi(0)}\mathcal{D}\left[\Psi^*,\Psi\right] 
		\text{e}^{-\int^\beta_0 d\tau\left(
				\sum_{ij}\Psi^*_i(t)\left[\delta_{ij}(\partial/\partial_\tau
		-\mu)+t_{ij}\right]\Psi_j(t)+\sum_{ijkl}U_{ijkl}\Psi^*_{i}(t)\Psi^*_{j}(t)\Psi_k(t)\Psi_l(t)\right)}
	\end{aligned}
\end{equation}
, where the shortened expression, $ \mathcal{D}\left[\Psi^*,\Psi\right] = \prod^{t=N-1 }_{t=1}
\left(\prod_id\Psi^*_i(t)d\Psi_i(t)\right) $ has been introduced. The sense of this expression is understood 
as time ordered integral of all possible functional over all possible paths. Algebraically, any transformation
for the Grassmann fields would also transform this integral measure by a factor of the Jacobian of the
transformation. The quantity in the exponent of the above expression is known as the Euclidean action 
or imaginary time action, 
$ S_{\text{E}}\left(\Psi^*,\Psi\right)= \int d\tau\left[\Psi^*\left(\partial/\partial_\tau\right)\Psi+H(\Psi^*,\Psi)\right]$
(with suppressed indices).

\subsection{Hubbard-Stratanovich Transformation}

The reformulation of the Fermionic problem in the language of functional integrals, or as it is more commonly referred
in the literature, the path integral, makes way for utilising various controlled approaches towards solving the problem. 
If the model that we are discussing for the moment did not have the four-fermion interaction term, it would have been 
readily soluble by diagonalising the hopping matrix $ \bm{t} $ in real space. If the underlying lattice of the model 
has translational symmetry then the dimensionality of the diagonalisation problem could be reduced by fourier transforming
the Grassmann variable. E.g. for a one-dimensional lattice without any additional sublattice degrees of freedom would 
transform the non-interacting Euclidean action into a simple diagonal form for the fourier modes, $ \Psi(t)_{\bm{k}} $. The 
differential operator can also be tackled by doing a temporal fourier transform, 
$ \Psi(t)_{\bm{k}}=\frac{1}{\sqrt{\beta}}\sum_{\omega_n}\Psi_{\bm{k}}(n)e^{-i\omega_n\tau}$. 
In this way the partion function can be reduced to multiple of Gaussian
integrals and can be exactly evaluated and can be found to agree  with results obtained from single particle quantum 
mechanical calculations.

However, the real benefit of the functional integral formalism emerges in the presence of many-body interactions when 
standard quantum-mechanical calculations break down. In the presence of weak interaction the exponential in the functional 
integral can be formally expanded in a Taylor series of interaction strength $ U $ and pertarbative results can be obtained. 
In principle the calulation in such a scheme would boil down to computing multiple orders of Gaussian integrals with the 
'free' two-body part of the action appearing as the quadratic exponent of the Gaussian integral. As a look from the expression
of the partition function would reveal, this is far from being a straightforward exercise as multiple mode of the 
Grassmann fields can combine in different ways in the expansion to give non-zero contribution. Wick's theorem (to be discussed
shortly) gives a formal way to collect these contributions. Feynman's diagram makes the job of applying Wick's theorem 
easier in a pictorial way. All interactions would not be weak and there are other non-perturbative methods like Wilsonian RG, where
the functional integral is carefully integrated out for the higher energy modes and thereby resulting into a functional 
integral of fields only involving the lower energy modes. In this procedure often the action changes dynamically and then 
assumes a simple ultimate form. A lot of interesting discovery of condensed matter physics in the past century came through 
this systamatic way of looking at a physical model in different energy scales. There are still other methods involving numerical 
computation of the functional integral using approximate techniques like the Monte-Carlo sampling. 


For this section we shall look at a particular method of recasting the interacting Fermionic functional integral problem
to a coupled but non-interacting Bosonic-Fermionic problem. This method, named as the Hubbard-Stratanovich transformation, 
remains an exact description of the original problem. However, it leads way to novel approximation schemes (as we 
shall see later) and can also be used to look at effective field theories at different energy scales. The transformation 
utilises basic identities of Grassmann and scalar Gaussian integrals which we list ahead of the formal H-S transformation of our
model problem.


\end{document}
